\documentclass[10pt,a4paper]{amsart}
\usepackage[margin=19mm]{geometry}
\usepackage{amsmath,amssymb,mathtools}
\usepackage[scr=rsfs,cal=euler]{mathalfa}
\usepackage{xfrac}
\usepackage[unicode,hidelinks,bookmarks=false]{hyperref}
\newcommand{\ie}{i.e.\ }
\newcommand{\cf}{cf.\ }
\newcommand{\resp}{respectively\ }
\newcommand{\Subsection}{\subsection}
\providecommand{\nobreakdash}{\nobreak\hbox{-}}
\setlength{\emergencystretch}{2em}
\multlinegap0pt
\usepackage{bm}
\usepackage{bbm}
\usepackage{dsfont}
\usepackage{xspace}
\usepackage{cancel}
\usepackage{mathtools}
\usepackage[shortlabels]{enumitem}
\usepackage{amsthm}
\makeatletter
\@ifundefined{theorem}{\newtheorem{theorem}{Theorem}}{}
\@ifundefined{lemma}{\newtheorem{lemma}[theorem]{Lemma}}{}
\makeatother
\makeatletter
\newcommand{\D}{\mathrm{dom}\,}
\newcommand{\DB}{\mathrm{dom}\,B }
\newcommand{\DT}{\mathrm{dom}\,T }
\newcommand{\adj}{^*}
\newcommand{\bhk}{B(\h,\kk)}
\newcommand{\bk}{B(\kk)}
\newcommand{\h}{\mathfrak{H}}
\newcommand{\half}{^{\frac{1}{2}}}
\newcommand{\halfltsb}{(\lambda  T^*B)^{\frac{1}{2}}}
\newcommand{\halfltsbs}{ (\lambda \,  \tstar B)^{\half}_s }
\newcommand{\halftsb}{( T^*B)^{\frac{1}{2}}}
\newcommand{\halftsbs}{ ( \tstar B)^{\half}_s }
\newcommand{\ifaf}{\Leftrightarrow}
\newcommand{\inclu}{\subseteq}
\newcommand{\kerr}{\mathrm{ker}\,}
\newcommand{\kk}{\mathfrak{K}}
\renewcommand{\labelenumi}{(\arabic{enumi})}
\newcommand{\lplusk}{B^{+}(\kk)}
\newcommand{\mul}{\mathrm{mul} \, }
\newcommand{\n}{\|}
\newcommand{\orth}{^\perp}
\expandafter\ifx\csname proof\endcsname\relax
\renewenvironment{proof}[1][\proofname]{%
   \par\pushQED{\qed}\normalfont%
   \topsep6\p@\@plus6\p@\relax
   \trivlist\item[\hskip\labelsep\bfseries#1\@addpunct{.}]%
   \ignorespaces
}{%
   \popQED\endtrivlist\@endpefalse
}
\fi
\newcommand{\ran}{\mathrm{ran}\,}
\newcommand{\ranbar}{\overline{\mathrm{ran}}\,}
\newcommand{\ranbarT}{\mathrm{\overline{ran}}\, T}
\newcommand{\rest}{\upharpoonright}
\newcommand{\ts}{T_s}
\newcommand{\tstar}{T^{*}}
\makeatother
\title[Product of nonnegative selfadjoint operators in unbounded se]{Product of nonnegative selfadjoint operators in unbounded settings}
\author{Yosra Barkaoui, Seppo Hassi}
\date{Source: arXiv:2507.14404v1 (CC BY 4.0)}
\begin{document}
\maketitle

\noindent\emph{Source and licence.} Product of nonnegative selfadjoint operators in unbounded settings. Version arXiv:2507.14404v1 (CC BY 4.0), distributed under the Creative Commons Attribution 4.0 International licence, as recorded in the arXiv licence metadata for this exact version and shown on the abstract page https://arxiv.org/abs/2507.14404v1. Full rights notice: licenses/arxiv-2507.14404-CC-BY-4.0.txt.\par\smallskip

\noindent\textit{Editorial scope.} Counted unit: the Theorem whose source label is \texttt{relationtheo} in the article (source0.tex lines 1290--1310, source label \texttt{relationtheo}), delivered together with the article's own complete original proof of it (source0.tex lines 1313--1426, one \texttt{proof} environment of 114 lines). No mathematics is written, repaired or re-derived: statement and proof are the article's own text line for line. Required context retained uncounted: lemprepthe (lines 769--787). Every definition, display and label of those lines is retained unchanged. Omitted from this package and not redistributed: 1--768, 788--1289, 1311--1312, 1427--2145. Nothing inside a retained excerpt is omitted or altered: every retained body is delivered exactly as the article's lines are written, so no omission receipt of any kind accompanies this package. Citations: the retained excerpts carry no citation command at all, so no bibliography entry of the article is transcribed and no reference is redistributed. Reference adaptation: none was needed and none was made; every reference inside the delivered unit resolves inside it, so no unresolved reference or citation is produced. Printed slips kept literal: no printed word, symbol, hypothesis or number of the counted statement or of its proof was altered. Layout and class: the article's own class is \texttt{amsart} with packages including amssymb, latexsym, euscript, graphicx, lmodern, pifont, ulem, amsmath, amscd, amsthm, amsfonts, color, babel, fontenc; the wrapper is the lane's pinned \texttt{amsart} editorial wrapper at 10pt on A4, which restates the theorem-like environments the retained text needs and adds the article's own macro definitions that the retained text needs (\texttt{\textbackslash D}, \texttt{\textbackslash DB}, \texttt{\textbackslash DT}, \texttt{\textbackslash adj}, \texttt{\textbackslash bhk}, \texttt{\textbackslash bk}, \texttt{\textbackslash h}, \texttt{\textbackslash half}, \texttt{\textbackslash halfltsb}, \texttt{\textbackslash halfltsbs}, \texttt{\textbackslash halftsb}, \texttt{\textbackslash halftsbs}, \texttt{\textbackslash ifaf}, \texttt{\textbackslash inclu}, \texttt{\textbackslash kerr}, \texttt{\textbackslash kk}, \texttt{\textbackslash labelenumi}, \texttt{\textbackslash lplusk}, \texttt{\textbackslash mul}, \texttt{\textbackslash n}, \texttt{\textbackslash orth}, \texttt{\textbackslash ran}, \texttt{\textbackslash ranbar}, \texttt{\textbackslash ranbarT}, \texttt{\textbackslash rest}, \texttt{\textbackslash ts}, \texttt{\textbackslash tstar}), with extra wrapper packages (bm, bbm, dsfont, xspace, cancel, mathtools, enumitem). The delivered numbering is the unit's own. Assets: no retained span contains a figure environment, a graphics-inclusion command, a PDF-page inclusion, a table or a TikZ picture; no image file is copied into this package, the package declares no asset dependency and no image file is redistributed with it. Licence: version arXiv:2507.14404v1 of this article is distributed under the Creative Commons Attribution 4.0 International licence, as recorded in the arXiv licence metadata for this exact version and shown on the abstract page https://arxiv.org/abs/2507.14404v1. A full rights notice is delivered at licenses/arxiv-2507.14404-CC-BY-4.0.txt. Characters: every retained excerpt is pure ASCII, so the delivered engine, XeLaTeX with its default Latin Modern text font, reports no missing character.
\begin{lemma} \label{lemprepthe}
Let $T,B:\h \rightarrow \kk $ be closed densely defined linear operators.  Then the following statements  are equivalent:
   \begin{enumerate}\def\labelenumi{\rm(\roman{enumi})}
     \item $Y B \inclu T$ has a solution $Y \in \bk;$
   \item   $T^* T  \leq c^2 B\adj B$  for some $0 \leq c \ (  =\n Y\n).$
      \end{enumerate}
     In this case, $Y$ can be selected such that $\ran Y \inclu \ranbarT$ and $\ker B\adj \inclu \ker Y.$ Furthermore,  if $\tstar B$ is selfadjoint then the following implication holds
            %-----------------------------------
      \begin{equation}\label{darkness1}
       Y B \inclu T  \text{ for some } Y \in \lplusk \Rightarrow T^* T  \leq  c_1  \tstar B \leq c_2 \, B\adj B,
      \end{equation}
      %-----------------------------------
       where $c_1,c_2 \geq 0.$
     In this case $\DB \inclu \D \halftsb \inclu  \DT$ and
       \begin{equation}\label{darkness2}
       \tstar B= B\adj Y B = B\adj T.
      \end{equation}
      %-----------------------------------
\end{lemma}

\begin{theorem}\label{relationtheo}
  Let $T, B: \h\rightarrow \kk$ be closed linear relations such that
$\mul B \subseteq \ker (\ts)\adj$ and   $\tstar B $ is selfadjoint.
Then, the following statements are equivalent  for $B_0:= B \rest \D \tstar B$ and for  some $\lambda \geq  0:$
 \begin{enumerate}\def\labelenumi{\rm(\roman{enumi})}
     \item $T^* T  \leq \lambda \, \tstar B ;$
     \item $  X \overline{B_0}\inclu T   $ has a solution $X \in \lplusk.$
          \end{enumerate}
 In this case, $X$ can be chosen such that $X\overline{B_0}  \inclu T_s$ and $ \ker (\ts)\adj \inclu \ker X$ with $\n X \n \leq \lambda.$ Moreover, in this case
      %--------------------------------
      \begin{equation}\label{hawnha}
        \tstar \overline{B_0} =  B_0 \adj X \overline{B_0}= B_0\adj T.
      \end{equation}
      %--------------------------------
 In particular, the following assertions are equivalent for some $\lambda \geq 0:$
   \begin{enumerate}[i)]\def\labelenumi{\rm(\roman{enumi})}\setcounter{enumi}{2}
     \item $T^* T  \leq \lambda \tstar \overline{B_0}$ and $\DT  \inclu \D \overline{B_0}$;
     \item $T=X\overline{B_0}\,\overset{.}{+}\, T_{\mul}$ has a solution $X \in \lplusk.$
      \end{enumerate}
      In this case, $X$ can be chosen such that $T_s=X\overline{B_s}$ and $\ker X =\ker (\ts)\adj.$
\end{theorem}

\begin{proof}
  Observe that  item $(i)$  is equivalent to $(\ts)\adj \ts \leq \lambda \halftsb_s \halftsb_s ,$ and hence the formula
     %
          $$
            \begin{array}{rcl}
                G: \ran   (\tstar B)^{\half}_s & \longrightarrow & \ran T_s \\
                 \halfltsbs  f               & \longmapsto &  T_s f,  \quad f \in \D (\tstar B)^{\frac{1}{2}}_s,
           \end{array}
           $$
   defines a contractive operator from $\ran (\tstar B)^{\half}_s$ into $\ran\ts,$ since
   $$
   \n G  \halfltsbs f \n^2 =
    \n  T_s f \|^2 \leq \n \halfltsbs  f \|^2
    $$
    for all $f \in \D   \halftsbs.$
   Moreover, $G $ can be extended to an operator $G_0 \in \bhk$ such that
      %_______________________________________
   \begin{eqnarray}\label{formulakerg}
     \ker( G_0 )\supseteq  ( \ran     \halftsbs )\orth
      = & \ker     \halftsbs  \oplus \mul\tstar B  \nonumber \\
      = & \hspace*{-0.5cm} \kerr   \tstar B \oplus \mul\tstar B
   \end{eqnarray}
   %_______________________________________
   and    $\ranbar G_0 \inclu \ranbar \ts $, which is equivalent to
   $\ran (G_0)\adj \inclu \ranbar(\tstar B)_s$ and
    \begin{equation}\label{correc}
      \kerr   G_0\adj \supseteq\kerr   (\ts)\adj=  \kerr   \tstar \oplus \mul T.
    \end{equation}
      %_______________________________________
   %_______________________________________
   Thus,
   \begin{equation}\label{otherprof}
     G_0   \halfltsbs \inclu T_s.
   \end{equation}
  As $T$  is closed, $\ts$ is also closed and $\ts \inclu T.$ Hence,
  \begin{equation*}\label{other5}
  \tstar \inclu (\ts)\adj \inclu      \halfltsb   G_0\adj
  \end{equation*}
  which implies that
  \begin{equation}\label{other55}
  \lambda \tstar B \inclu \halfltsb \lambda G_0\adj B.
 \end{equation}
  By assumption $\mul B \inclu \ker (\ts)\adj$, and hence $\mul \tstar B = \mul \tstar T = \mul \tstar.$    On the other hand,   $\mul B \inclu \ker (\ts)\adj \inclu \ker   G_0\adj$ (see \eqref{correc}), so
          %
       \begin{equation} \label{other0}
         G_0\adj B   = G_0\adj  (B_s  \overset{.}{+} B_{\mul})
                          =    G_0\adj  B_s.
        \end{equation}
            This yields by \eqref{other55} that
            %
            %__________________________
                 \begin{equation} \label{other}
                  (\lambda  \tstar B)_s \inclu    \halfltsbs G_0\adj B_s.
                   \end{equation}
                   %
              %_______________________________________
        Multiplying \eqref{other} from the left by the Moore-Penrose inverse $  (\tstar B)_s^{(-\frac{1}{2})}$  gives
            %-------------------------------
                 \begin{align*}
                           Q_s (I \rest \D \halftsbs)   \halfltsbs & \inclu   Q_s I \rest \D \halftsbs   \lambda  G_0\adj B_s\\
                                                           & \inclu  \lambda Q_s    G_0\adj B_s = \lambda G_0\adj B_s,
                \end{align*}
               %-------------------------------
               where $Q_s$ is the orthogonal  projection onto $\ranbar (\tstar B)_s.$
               Consequently,
              \begin{equation*}
                \halfltsbs \rest \D \tstar B \inclu \lambda   G_0\adj B_s,
              \end{equation*}
              and hence
              $
                  \halfltsbs| \D \tstar B  = \lambda  G_0\adj B_s \rest \D \tstar B.
              $
             Since $\D \tstar B $ is a core for $ \halftsbs$,   one gets
               \begin{equation*}
               \halfltsbs =   \lambda \overline{G_0\adj B_s \rest \D \tstar B} \supseteq \lambda  G_0\adj \overline{ B_s \rest \D \tstar B}.
              \end{equation*}
              Together with \eqref{otherprof} and \eqref{other0}  this implies that
             \begin{equation}\label{ghayth}
              \lambda G_0 G_0\adj \overline{B_0}= \lambda  G_0 G_0\adj \overline{ B_s \rest \D \tstar B} \inclu  G_0  \halfltsbs \inclu \ts
               \end{equation}
            and, in particular,
              $$
                \lambda G_0 G_0\adj \overline{ B_0} \inclu  T.
                $$
         This proves (ii) for $X=  \lambda  G_0 G_0\adj \in \lplusk.$

                 \hspace*{0.4cm}  The inclusion $\kerr T\adj \inclu \ker (\ts)\adj  \inclu \kerr X$ follows  from \eqref{correc}  by fact that  $\ker (G_0)\adj= \ker X$.


               For the reverse implication $\rm(ii) \Rightarrow (i)$, observe that
   $$
   \tstar \overline{B_0} \inclu B_0\adj X \overline{B_0} \inclu (X\half \overline{B_0})\adj \overline{X\half \overline{B_0}},
   $$
    and since  $ \tstar \overline{B_0} =\tstar B $ is selfadjoint also   $B_0\adj X \overline{B_0}$    is   selfadjoint. Thus,
   \begin{equation}\label{9rib}
      \tstar \overline{B_0}= B_0\adj X \overline{B_0}.
  \end{equation}
  Now, $X\half \overline{X\half \overline{B_0} } \inclu \overline{X\overline{B_0}} \inclu T$ and the same argument that was used in the proof of  Lemma \ref{lemprepthe} shows that for  $\lambda = \n X \n $  one has
  $$
  \tstar T \leq \lambda  (X\half B)\adj \overline{X\half \overline{B_0}} = \lambda  B_0\adj X\overline{B_0}  = \lambda \tstar \overline{B_0},
  $$
  and (i) is proved.


 %-----------------
To complete the proof of  \eqref{hawnha}  observe that
 $$
    B_0\adj T \inclu (\tstar \overline{B_0})\adj= \tstar \overline{B_0} = B_0\adj X \overline{B_0} \inclu B_0\adj T,
 $$
  %------------------
  and hence $B_0\adj T= B_0\adj X \overline{B_0} =\tstar \overline{B_0}.$\\
   \hspace*{0.4cm}  For the proof of the equivalence $\rm(iii) \ifaf (iv)$, it suffices to observe that item (iv) is   equivalent to  $
  X \overline{B_0} \inclu \ts  $ and $ \DT = \D \overline{B_0},$ and conclude the result from the first part of the proof. In this particular case, one easily sees that $\ts =X \overline{B_0} $ and hence $(\ts)\adj = B_0\adj X$, which leads to  $\ker X \inclu \ker  (\ts)\adj \inclu \ker X.$
\end{proof}

\end{document}
